\documentclass[aspectratio=169,10pt]{beamer}
\usetheme{default}
\setbeamertemplate{navigation symbols}{}
\usepackage{graphicx}\usepackage[T1]{fontenc}\usepackage{amsmath,amssymb}
\usepackage{listings}
\definecolor{ncsured}{HTML}{CC0000}\definecolor{ncsudk}{HTML}{990000}
\definecolor{ncsugray}{HTML}{6D6D6D}\definecolor{ncsuteal}{HTML}{008473}
\definecolor{codebg}{HTML}{F5F5F5}
\setbeamercolor{structure}{fg=ncsured}
\setbeamercolor{frametitle}{fg=white,bg=ncsured}
\setbeamercolor{title}{fg=ncsudk}
\setbeamercolor{itemize item}{fg=ncsured}
\setbeamercolor{itemize subitem}{fg=ncsured}
\setbeamerfont{frametitle}{series=\bfseries}
\setbeamertemplate{frametitle}{\nointerlineskip\vskip2pt%
  \begin{beamercolorbox}[wd=\paperwidth,ht=2.4ex,dp=1ex,leftskip=.3cm]{frametitle}%
  \small\bfseries\insertframetitle\end{beamercolorbox}}
\setbeamertemplate{footline}{\hbox{\begin{beamercolorbox}[wd=\paperwidth,ht=2.2ex,dp=1ex,leftskip=.3cm,rightskip=.3cm]{}%
  \tiny Optimal Market Making for Cryptocurrency Options \hfill Week 7 \hfill \insertframenumber\end{beamercolorbox}}}
\lstset{basicstyle=\ttfamily\scriptsize,backgroundcolor=\color{codebg},
  keywordstyle=\color{ncsured},commentstyle=\color{ncsugray},
  stringstyle=\color{ncsuteal},breaklines=true,columns=fullflexible,
  showstringspaces=false,frame=single,rulecolor=\color{ncsugray}}

\title{Week 7: Fast Computation: Prices, Greeks, Implied Vol, Surfaces}
\subtitle{The latency budget, in C++ via pybind11 -- Milestone 2}
\date{North Carolina State University \,\textbar\, Summer 2026}
\titlegraphic{
\includegraphics[width=4.2cm]{/Users/dendisuhubdy/Github/ncsu/assets/ncsu_logo.png}}
\begin{document}
\begin{frame}[plain]\titlepage\begin{center}\tiny Industry Mentor: Dendi Suhubdy (Bitwyre)  |  Guest: Fenni Kang (Coincall)\end{center}\end{frame}
\begin{frame}{Learning Objectives}
\begin{itemize}
  \item Implement COS / Carr-Madan pricers and robust implied-vol in C++
  \item Expose the C++ kernels to Python with pybind11
  \item Validate Greeks against PyTorch autograd; hit the 5 ms budget
\end{itemize}
\end{frame}

\begin{frame}{Lecture 1: CF Pricing \& Implied Vol}
\begin{itemize}
  \item Characteristic functions: Black-Scholes and Heston
  \item Fang-Oosterlee COS method (the workhorse)
  \item Robust implied-vol inversion (good guess + safeguarded Newton)
\end{itemize}
\end{frame}

\begin{frame}{Lecture 2: Binding C++ with pybind11}
\begin{itemize}
  \item Structure a pybind11 module; build with CMake
  \item Pass NumPy/PyTorch tensors with no copy (buffer protocol)
  \item Release the GIL around the C++ hot loop
  \item PyTorch autograd is the REFERENCE for Greeks; C++ is the speed
\end{itemize}
\end{frame}

\begin{frame}{Illustration}
\begin{center}\includegraphics[width=0.86\textwidth]{/Users/dendisuhubdy/Github/ncsu/weekly_presentations/figures/week07.png}\end{center}
\end{frame}

\section*{Derivation}
\begin{frame}{Derivation \& Rationale}
\[C(K)=e^{-rT}\,\tfrac{1}{\pi}\!\int_0^\infty \mathrm{Re}\!\big[e^{-iu\ln K}\hat{\phi}(u)\big]\,du\]

\small\textbf{\color{ncsured}Why:} Once we leave Black-Scholes (e.g. Heston), there is no closed-form price, but the characteristic function $\phi$ IS known. Pricing becomes a single Fourier integral -- this is why CF methods exist and why they are fast.
\end{frame}

\begin{frame}{Derivation \& Rationale}
\[C \approx \sum_{n=0}^{N-1}{}' \mathrm{Re}\!\Big[\phi\!\big(\tfrac{n\pi}{b-a}\big)e^{-i n\pi\frac{a}{b-a}}\Big]V_n\]

\small\textbf{\color{ncsured}Why:} The Fang-Oosterlee COS method replaces the integral with a cosine series that converges geometrically for smooth densities. We truncate $[a,b]$ from the cumulants, not by guessing -- that is what makes it reach $10^{-4}$ accuracy in microseconds.
\end{frame}

\begin{frame}{Getting the Data from Coincall}
\begin{itemize}
  \item Inputs are the calibrated surface from Week 6 plus reference prices:
  \begin{itemize}\item \texttt{\small coincall\_btc\_options\_chain\_snapshots\_2025Q4.parquet}\end{itemize}
  \begin{itemize}\item \texttt{\small coincall\_calibrated\_svi\_surfaces\_2025Q4.parquet (baseline)}\end{itemize}
  \begin{itemize}\item \texttt{\small coincall\_reference\_prices\_2025Q4.parquet (accuracy target, 1e-4)}\end{itemize}
  \item Live mark/IV/Greeks for spot checks: GET /open/option/detail/v1/\{symbol\}
\end{itemize}
\vspace{2pt}\tiny\textit{Cross-check exact paths at https://docs.coincall.com/}
\end{frame}

\begin{frame}[fragile]{Code --- fastmm pybind11 module}
\begin{lstlisting}
// fastmm.cpp  (build with CMake + pybind11; import as `fastmm`)
#include <pybind11/pybind11.h>
#include <pybind11/numpy.h>
namespace py = pybind11;

double cos_price(double S, double K, double T, double r, double sigma);  // COS pricer

py::array_t<double> price_surface(py::array_t<double> K, double S,
                                  double T, double r, double sigma) {
    auto k = K.unchecked<1>(); auto out = py::array_t<double>(k.shape(0));
    auto o = out.mutable_unchecked<1>();
    { py::gil_scoped_release nogil;                 // release GIL for the hot loop
      for (py::ssize_t i = 0; i < k.shape(0); ++i)
          o(i) = cos_price(S, k(i), T, r, sigma); }
    return out;                                     // zero-copy back to NumPy
}
PYBIND11_MODULE(fastmm, m) { m.def("price_surface", &price_surface); }
\end{lstlisting}
\end{frame}
\begin{frame}{Glossary --- Key Terms}
\footnotesize\begin{description}
\item[\color{ncsured}Characteristic function] the Fourier transform of the log-return density
\item[\color{ncsured}COS method] Fourier-cosine series option pricing (Fang-Oosterlee)
\item[\color{ncsured}Carr-Madan] FFT-based option pricing via a damped payoff
\item[\color{ncsured}Implied-vol inversion] solving for the IV that reproduces a price (safeguarded Newton)
\item[\color{ncsured}pybind11] header-only C++ \textless{}-\textgreater{} Python binding
\item[\color{ncsured}GIL] Python's Global Interpreter Lock; released around the C++ hot loop
\item[\color{ncsured}Latency budget] the maximum allowed time (sub-5 ms full-surface revaluation)
\end{description}
\end{frame}

\begin{frame}{Reading}
\begin{itemize}
  \item Carr-Madan (1999); Fang-Oosterlee (2008); Jaeckel (2015); pybind11 docs
\end{itemize}
\end{frame}

\begin{frame}{References}
\scriptsize\begin{itemize}\setlength{\itemsep}{3pt}
  \item Carr, P., \& Madan, D. (1999). Option valuation using the fast Fourier transform. J. Computational Finance, 2(4), 61-73.
  \item Fang, F., \& Oosterlee, C. W. (2008). A novel pricing method for European options based on Fourier-cosine series expansions. SIAM J. Sci. Comput., 31(2), 826-848.
  \item Jackel, P. (2015). Let's be rational. Wilmott Magazine.
  \item Jakob, W., Rhinelander, J., \& Moldovan, D. (2017). pybind11. https://github.com/pybind/pybind11
  \item Paszke, A., et al. (2019). PyTorch: an imperative style, high-performance deep learning library. NeurIPS.
\end{itemize}
\end{frame}

\end{document}